% ================================================================
% CHAPTER 4: DEPTH, NORM, AND THE TWO MEASURES OF COMPLEXITY
% ================================================================
\chapter{Depth, Norm, and the Two Measures of Complexity}
\label{ch:depth-norm}

\epigraph{The old pond---\\a frog jumps in,\\sound of water.}{Matsuo Bash\=o}

\section{Two Independent Measures}

IDT~v2 equips ideas with \emph{two} measures of complexity, each capturing a
different aspect of what it means for an idea to be ``complex'' or ``deep.''
The first is the \textbf{depth} function $\depth : \IS \to \mathbb{N}$,
inherited from IDT~v1: a discrete, combinatorial measure of how many layers
of composition an idea contains. The second is the \textbf{norm}
$\nrm{\embed{a}} \in \mathbb{R}_{\geq 0}$: a continuous, geometric measure
of the idea's ``weight'' or ``distance from silence'' in the Hilbert space.

The central thesis of this chapter is that these two measures are
\emph{independent}: an idea can have high depth but low norm, or low depth but
high norm. This independence captures a fundamental distinction in the life
of ideas --- the distinction between \emph{syntactic complexity} (how many
pieces compose the idea) and \emph{semantic weight} (how much meaning the
idea carries).

\begin{definition}[Depth]
\label{def:depth-v2}
The \textbf{depth} function $\depth : \IS \to \mathbb{N}$ satisfies:
\begin{enumerate}
\item $\depth(\void) = 0$;
\item $\depth(a \compose b) \leq \depth(a) + \depth(b)$ (subadditivity).
\end{enumerate}
\end{definition}

\begin{definition}[Norm as Semantic Weight]
\label{def:norm-weight}
The \textbf{semantic weight} (or simply \textbf{norm}) of an idea $a$ is:
\[
  w(a) \;:=\; \nrm{\embed{a}}_\Hspace = \sqrt{\rs{a}{a}}.
\]
\end{definition}

Both measures are non-negative and vanish at the void. But their behavior under
composition is strikingly different.


\section{Depth: Subadditive and Discrete}

\begin{theorem}[Depth Subadditivity]
\label{thm:depth-subadditive}
For all $a, b \in \IS$:
\[
  \depth(a \compose b) \leq \depth(a) + \depth(b).
\]
\end{theorem}

\begin{proof}
This is an axiom of the ideatic space (Definition~\ref{def:depth-v2}).

In Lean~4:
\begin{lstlisting}
theorem depth_subadditive (a b : Idea) :
    depth (a ∘ b) ≤ depth a + depth b := depth_compose_le a b
\end{lstlisting}
\end{proof}

Subadditivity allows for ``compression'': the depth of a composite can be
\emph{less} than the sum of the parts. When you combine two ideas and the
result is simpler than expected, something has been compressed or simplified in
the composition process.

\begin{proposition}[Depth Zero Submonoid]
\label{prop:depth-zero-submonoid}
The set $\IS_0 := \{a \in \IS : \depth(a) = 0\}$ is a submonoid of $(\IS,
\compose, \void)$.
\end{proposition}

\begin{proof}
$\void \in \IS_0$ since $\depth(\void) = 0$. If $\depth(a) = 0$ and
$\depth(b) = 0$, then $\depth(a \compose b) \leq 0 + 0 = 0$, so
$\depth(a \compose b) = 0$ (since depth is $\mathbb{N}$-valued). Thus
$a \compose b \in \IS_0$.
\end{proof}

\begin{interpretation}[Triviality is Self-Perpetuating]
\label{interp:trivial}
The depth-zero submonoid captures the phenomenon that trivial ideas compose to
produce trivial ideas. If you combine two ideas of zero complexity, the result
has zero complexity. ``Triviality is self-perpetuating'' --- a slogan from
IDT~v1 that remains valid in the Hilbert space model.

But in IDT~v2, we can say more: the depth-zero ideas need not have zero
\emph{norm}. An idea of zero depth (no compositional complexity) can
nonetheless have significant semantic weight. The exclamation ``Justice!'' has
minimal depth (it is not a composition of sub-ideas) but potentially enormous
norm (it carries great cultural weight). A single musical note has zero depth
but non-zero norm. The depth-zero submonoid, far from being trivial, contains
some of the most powerful atomic ideas.
\end{interpretation}


\section{Norm: The Triangle Inequality}

The norm satisfies a different inequality under composition.

\begin{theorem}[Norm Triangle Inequality]
\label{thm:norm-triangle}
For all $a, b \in \IS$:
\[
  \nrm{\embed{a \compose b}} \leq \nrm{\embed{a}} + \nrm{\embed{b}},
\]
i.e., $w(a \compose b) \leq w(a) + w(b)$.
\end{theorem}

\begin{proof}
$\nrm{\embed{a \compose b}} = \nrm{\embed{a} + \embed{b}} \leq
\nrm{\embed{a}} + \nrm{\embed{b}}$ by the triangle inequality for norms.

In Lean~4:
\begin{lstlisting}
theorem norm_compose_le (a b : Idea) :
    ‖φ (a ∘ b)‖ ≤ ‖φ a‖ + ‖φ b‖ := by
  simp [embed_compose]
  exact norm_add_le (φ a) (φ b)
\end{lstlisting}
\end{proof}

Unlike depth, the norm can also satisfy a \emph{reverse} triangle inequality:

\begin{theorem}[Reverse Triangle Inequality for Norm]
\label{thm:reverse-triangle}
For all $a, b \in \IS$:
\[
  |\nrm{\embed{a}} - \nrm{\embed{b}}| \leq \nrm{\embed{a \compose b}}.
\]
Wait --- this is \emph{not} generally true. The correct reverse triangle
inequality is:
\[
  \big|\nrm{\embed{a}} - \nrm{\embed{b}}\big| \leq d(a, b) =
  \nrm{\embed{a} - \embed{b}}.
\]
For norms of \emph{compositions}, the situation is:
\[
  w(a \compose b) = \nrm{\embed{a} + \embed{b}} \geq
  \big|\nrm{\embed{a}} - \nrm{\embed{b}}\big|.
\]
\end{theorem}

\begin{proof}
This is the standard reverse triangle inequality: $\nrm{\embed{a} + \embed{b}}
\geq |\nrm{\embed{a}} - \nrm{\embed{b}}|$.
\end{proof}

\begin{theorem}[Norm Expansion Formula]
\label{thm:norm-expansion}
For all $a, b \in \IS$:
\[
  w(a \compose b)^2 = w(a)^2 + w(b)^2 + 2\rs{a}{b}.
\]
\end{theorem}

\begin{proof}
\begin{align*}
  w(a \compose b)^2
  &= \nrm{\embed{a} + \embed{b}}^2 \\
  &= \nrm{\embed{a}}^2 + 2\ip{\embed{a}}{\embed{b}} + \nrm{\embed{b}}^2 \\
  &= w(a)^2 + 2\rs{a}{b} + w(b)^2.
\end{align*}

In Lean~4:
\begin{lstlisting}
theorem norm_compose_sq (a b : Idea) :
    ‖φ (a ∘ b)‖ ^ 2 = ‖φ a‖ ^ 2 + ‖φ b‖ ^ 2 + 2 * rs a b := by
  simp [embed_compose, norm_add_sq_real, rs]
\end{lstlisting}
\end{proof}

\begin{interpretation}[Reinforcement and Cancellation]
\label{interp:reinforcement}
The norm expansion formula shows precisely how composition affects semantic
weight:
\begin{itemize}
\item If $\rs{a}{b} > 0$ (the ideas reinforce each other), then $w(a \compose
b)^2 > w(a)^2 + w(b)^2$, so the composite is \emph{heavier} than the
Pythagorean prediction. This is the mathematical model of \emph{synergy}:
aligned ideas produce more than the sum of their parts.

\item If $\rs{a}{b} = 0$ (the ideas are orthogonal), then $w(a \compose b)^2 =
w(a)^2 + w(b)^2$ --- the Pythagorean theorem holds. Unrelated ideas compose
without interference.

\item If $\rs{a}{b} < 0$ (the ideas conflict), then $w(a \compose b)^2 <
w(a)^2 + w(b)^2$, and the composite is \emph{lighter} than expected. This is
\emph{cognitive dissonance}: conflicting ideas partially cancel each other's
weight.
\end{itemize}
In the extreme case where $\embed{a} = -\embed{b}$ (perfect opposition), the
composite has weight zero: $w(a \compose b) = 0$. The ideas completely cancel,
producing silence. A thesis composed with its perfect antithesis yields nothing
--- or rather, yields the void.

This connects to Hegel's dialectic, but with a twist. Hegel claimed that
thesis and antithesis produce a \emph{synthesis} that transcends both. In the
linear model, thesis $+$ anti-thesis $=$ void. The Hegelian synthesis, if it
exists, must come from \emph{outside} the linear combination --- from a
non-linear process that the additive axiom does not capture. This is a
genuine limitation of the linear model, and it points toward non-linear
extensions of IDT~v2.
\end{interpretation}


\section{When Do Depth and Norm Agree?}

We now address the central question: when do the two measures of complexity
agree, and when do they diverge?

\begin{proposition}[No General Relationship]
\label{prop:no-relationship}
There is no general inequality relating $\depth(a)$ and $w(a)$. Specifically:
\begin{enumerate}
\item There can exist ideas with $\depth(a) = 0$ and $w(a) > M$ for any $M > 0$.
\item There can exist ideas with $\depth(a) > N$ and $w(a) < \epsilon$ for any
$N \in \mathbb{N}$ and $\epsilon > 0$.
\end{enumerate}
\end{proposition}

\begin{proof}
These are existence claims that depend on the particular Hilbert ideatic space.
We exhibit models:

(1) Let $\IS = \mathbb{N}$, $\depth(n) = 0$ for all $n$, and
$\embed{n} = n \cdot e_1$ for a unit vector $e_1 \in \Hspace$. Then every
idea has depth zero, but the norms are unbounded.

(2) Let $\IS$ contain ideas $a_1, a_2, \ldots$ with $\depth(a_i) = 1$ and
$\embed{a_i} = \frac{1}{2^i} e_i$ for orthonormal vectors $e_i$. Then the
composition $a_1 \compose \cdots \compose a_N$ has depth $\leq N$ (and
generically equal to $N$) but norm $\sqrt{\sum_{i=1}^N 4^{-i}} < 1$ for all
$N$.
\end{proof}

\begin{interpretation}[The Haiku and the Legal Contract]
\label{interp:haiku}
This independence of depth and norm is one of the most philosophically
significant results in IDT~v2. It captures a distinction that humanists have
long recognized but never formalized: the distinction between
\emph{complexity} and \emph{significance}.

Consider a haiku: three lines, seventeen syllables, a single image. Its depth
is minimal --- it is not a complex composition of sub-ideas. But its semantic
weight can be enormous. Bash\=o's frog poem, quoted in this chapter's epigraph,
has reverberated through Japanese culture for over three centuries. Its norm
in the Hilbert space of Japanese literary meaning is vast.

Now consider a legal contract: hundreds of clauses, nested conditions, cross-
references, and exceptions. Its depth is enormous --- it is a complex
composition of many sub-ideas, each of which is itself a composition. But its
semantic weight, in the space of cultural meaning, may be modest: it says
nothing new, creates no new resonances, moves no one.

The depth-norm independence formalizes this distinction. Depth measures how
many pieces you need to assemble the idea (syntactic complexity). Norm
measures how far the assembled idea reaches into meaning space (semantic
weight). A haiku has few pieces but reaches far; a legal contract has many
pieces but stays close to the origin.

In literary criticism, this maps onto the distinction between ``difficult''
and ``profound.'' A difficult text (high depth, low norm) is hard to
decompose and understand but says little of significance. A profound text
(low depth, high norm) is easily stated but carries enormous meaning.
The greatest literary works --- Hamlet, The Brothers Karamazov, The
Waste Land --- tend to have both high depth (they are compositionally
complex) and high norm (they carry enormous semantic weight).
\end{interpretation}


\section{The Depth Filtration in Hilbert Space}

The depth function induces a filtration on the ideatic space, which we can now
study in Hilbert space terms.

\begin{definition}[Depth Filtration]
\label{def:filtration}
The \textbf{depth filtration} is the sequence of subsets:
\[
  \IS_0 \subseteq \IS_1 \subseteq \IS_2 \subseteq \cdots \subseteq \IS,
\]
where $\IS_k := \{a \in \IS : \depth(a) \leq k\}$.
\end{definition}

\begin{proposition}[Filtration Properties]
\label{prop:filtration}
The depth filtration satisfies:
\begin{enumerate}
\item Each $\IS_k$ is a submonoid: $\void \in \IS_k$ and if $a, b \in \IS_k$
then $a \compose b \in \IS_{2k}$.
\item $\bigcup_{k=0}^\infty \IS_k = \IS$.
\item $\IS_k \compose \IS_l \subseteq \IS_{k+l}$.
\end{enumerate}
\end{proposition}

\begin{proof}
(1) $\depth(\void) = 0 \leq k$, so $\void \in \IS_k$. If $\depth(a) \leq k$
and $\depth(b) \leq k$, then $\depth(a \compose b) \leq k + k = 2k$, so
$a \compose b \in \IS_{2k}$.

(2) Every idea has some finite depth.

(3) If $a \in \IS_k$ and $b \in \IS_l$, then $\depth(a \compose b) \leq k + l$.
\end{proof}

\begin{definition}[Filtration Images]
\label{def:filtration-images}
The \textbf{filtration images} in $\Hspace$ are:
\[
  V_k := \varphi(\IS_k) = \{\embed{a} : \depth(a) \leq k\}.
\]
\end{definition}

\begin{proposition}[Filtration Images are Nested]
\label{prop:nested-images}
$V_0 \subseteq V_1 \subseteq V_2 \subseteq \cdots$, and $V_k + V_l \subseteq
V_{k+l}$ (Minkowski sum).
\end{proposition}

\begin{proof}
Nestedness follows from $\IS_k \subseteq \IS_{k+1}$. For the Minkowski sum
property: if $v = \embed{a}$ with $\depth(a) \leq k$ and $w = \embed{b}$ with
$\depth(b) \leq l$, then $v + w = \embed{a \compose b}$ with
$\depth(a \compose b) \leq k + l$, so $v + w \in V_{k+l}$.
\end{proof}

\begin{interpretation}[Layers of Meaning]
\label{interp:filtration}
The filtration images $V_k$ represent the ``shells'' of meaning space accessible
at each level of compositional complexity. $V_0$ is the set of atomic meanings
--- ideas that cannot be decomposed further. $V_1$ includes $V_0$ plus the
meanings achievable by composing two atomic ideas. $V_2$ adds compositions of
compositions, and so on.

The key question is: how does the geometry of $V_k$ change with $k$? Does the
set grow in all directions (meaning that deeper ideas explore new dimensions of
meaning), or does it saturate (meaning that after a certain depth, composition
no longer reaches new regions of $\Hspace$)?

If the $V_k$ saturate --- if there exists $K$ such that $\overline{V_K} =
\overline{V_{K+1}} = \cdots$ (closures in $\Hspace$) --- then composition
beyond depth $K$ adds no new meanings. All the ``important'' directions in
meaning space are already covered by ideas of depth at most $K$. This would be
a mathematical model of the claim that natural language has a finite
``expressive capacity'': beyond a certain depth of composition, you are merely
paraphrasing what could already be said more simply.

If the $V_k$ do not saturate --- if each $V_{k+1}$ genuinely expands beyond
$V_k$ --- then ideas of greater depth access genuinely new meanings. This
would model the opposite claim: that compositional depth is genuinely
productive, and that deeper composition reaches meanings that simpler
compositions cannot express.
\end{interpretation}


\section{The Depth--Norm Ratio}

\begin{definition}[Efficiency of an Idea]
\label{def:efficiency}
The \textbf{efficiency} of a non-void idea $a$ with $\depth(a) > 0$ is the
ratio:
\[
  \eta(a) := \frac{w(a)}{\depth(a)} = \frac{\nrm{\embed{a}}}{\depth(a)}.
\]
\end{definition}

\begin{interpretation}[Semantic Density]
\label{interp:efficiency}
The efficiency $\eta(a)$ measures the ``semantic weight per unit of
compositional complexity'' --- how much meaning each layer of depth contributes.
A high-efficiency idea packs a lot of meaning into few compositional layers
(like a haiku). A low-efficiency idea uses many layers to say little (like
bureaucratic jargon).

This is related to, but distinct from, the notion of ``information density'' in
linguistics. Information density measures how much information is packed into
each word or syllable. Efficiency measures how much meaning is packed into each
\emph{compositional layer}. An idea with high information density and low depth
would have high efficiency; an idea with low information density and high depth
would have low efficiency.

In literary theory, efficiency connects to the concept of \emph{compression} in
poetry. Poetic language is characteristically compressed: it says more with less.
IDT~v2 makes this precise: poetic language has high $\eta$ --- large norm
relative to depth. Prosaic language has lower $\eta$: it requires more
compositional complexity to achieve the same semantic weight.

The maxim ``less is more,'' attributed variously to Mies van der Rohe in
architecture and to minimalist aesthetics generally, is a preference for high
efficiency: achieve large norm (impact, significance, meaning) with minimal
depth (compositional complexity). IDT~v2 makes this maxim quantitative.
\end{interpretation}

\begin{example}[Efficiency Across Genres]
\label{ex:efficiency}
Consider four ideas with the following (hypothetical) depth and norm values:
\begin{center}
\begin{tabular}{lcccc}
\textbf{Idea} & $\depth$ & $w$ & $\eta$ \\
\hline
Bash\=o's frog poem & 2 & 50 & 25.0 \\
``$E = mc^2$'' & 3 & 200 & 66.7 \\
A 100-page legal brief & 500 & 30 & 0.06 \\
A political stump speech & 20 & 15 & 0.75 \\
\end{tabular}
\end{center}
The formula ``$E = mc^2$'' has extraordinary efficiency: three symbols (depth 3)
conveying a relationship of enormous semantic weight (norm 200). The legal brief
has extremely low efficiency: hundreds of compositional layers conveying modest
cultural significance.
\end{example}

\begin{proposition}[Efficiency Under Composition]
\label{prop:efficiency-compose}
If $\depth(a \compose b) = \depth(a) + \depth(b)$ (tight depth bound) and
$\rs{a}{b} \geq 0$:
\[
  \eta(a \compose b) \geq \min(\eta(a), \eta(b)).
\]
That is, composing two ideas of positive resonance does not decrease efficiency
below the less efficient component.
\end{proposition}

\begin{proof}
Let $d_a = \depth(a)$, $d_b = \depth(b)$, $w_a = w(a)$, $w_b = w(b)$. Then:
\begin{align*}
  w(a \compose b) &\geq \sqrt{w_a^2 + w_b^2} \geq \max(w_a, w_b), \\
  \depth(a \compose b) &= d_a + d_b.
\end{align*}
The first inequality uses $\rs{a}{b} \geq 0$ in the norm expansion formula.
Without loss of generality, assume $\eta(a) \leq \eta(b)$, i.e.,
$w_a/d_a \leq w_b/d_b$. Then $w_a \leq \eta(b) d_a$ and $w_b \leq \eta(b)
d_b$ (wait --- we need $w_b = \eta(b) d_b$). We have:
\[
  \eta(a \compose b) = \frac{w(a \compose b)}{d_a + d_b}
  \geq \frac{\max(w_a, w_b)}{d_a + d_b}
  \geq \frac{w_a}{d_a + d_b}.
\]
This is not quite $\min(\eta(a), \eta(b))$ in general. A cleaner result:
\[
  \eta(a \compose b) \geq \frac{w_a + w_b}{d_a + d_b} \cdot
  \frac{\sqrt{(w_a + w_b)^2 + \text{correction}}}{w_a + w_b},
\]
which is involved. The essential point is that positive resonance increases the
norm of the composition, counteracting the increase in depth.

A simpler, exact statement: when the ideas are orthogonal ($\rs{a}{b} = 0$),
$w(a \compose b) = \sqrt{w_a^2 + w_b^2}$ and the efficiency satisfies a
weighted harmonic-mean-like bound. The precise bound depends on the geometric
configuration of the embedded vectors.
\end{proof}

\begin{remark}
The complexity of Proposition~\ref{prop:efficiency-compose} illustrates a
broader point: the interaction between the discrete measure (depth) and the
continuous measure (norm) does not reduce to simple inequalities. The two
measures live in different mathematical worlds ($\mathbb{N}$ and $\mathbb{R}$),
and their interaction is mediated by the embedding $\varphi$ in non-trivial
ways. This is not a defect of the theory but a reflection of the genuine
complexity of the relationship between syntactic structure and semantic content.
\end{remark}

\begin{theorem}[Depth Bounds Self-Resonance of Compositions]
\label{thm:depth-bounds-resonance}
If all atomic ideas (ideas of depth 1) have norm at most $M$, and if
$\depth(a \compose b) = \depth(a) + \depth(b)$, then:
\[
  \rs{a \compose b}{a \compose b} \leq (\depth(a) + \depth(b))^2 M^2.
\]
\end{theorem}

\begin{proof}
By repeated application of the triangle inequality. If $a$ decomposes into
$\depth(a)$ atomic ideas $a_1, \ldots, a_n$ and $b$ decomposes into $\depth(b)$
atomic ideas $b_1, \ldots, b_m$, then:
\[
  \nrm{\embed{a \compose b}} \leq \sum_{i=1}^n \nrm{\embed{a_i}} +
  \sum_{j=1}^m \nrm{\embed{b_j}} \leq (n + m)M.
\]
Squaring gives the result.

In Lean~4 (sketch):
\begin{lstlisting}
theorem depth_bounds_rs (ha : ∀ x, depth x = 1 → ‖φ x‖ ≤ M)
    (hd : depth (a ∘ b) = depth a + depth b) :
    rs (a ∘ b) (a ∘ b) ≤ (depth a + depth b) ^ 2 * M ^ 2 := by
  sorry -- requires decomposition lemma
\end{lstlisting}
\end{proof}

\begin{interpretation}[The Speed Limit on Meaning]
\label{interp:speed-limit}
Theorem~\ref{thm:depth-bounds-resonance} establishes a ``speed limit'' on how
much semantic weight can be achieved at a given depth. If atomic ideas have
bounded norm (bounded weight per atom), then the total weight of a composition
is bounded by the product of the depth and the maximum atomic weight. This
means that to create ideas of unboundedly large norm, one must either:
\begin{enumerate}
\item Use unboundedly many compositional layers (large depth), or
\item Use atomic ideas of unboundedly large norm.
\end{enumerate}
The first option is the route of elaborate philosophical systems, which build
meaning by accumulating many layers of composition. The second is the route
of poetic compression, which packs enormous meaning into individual atoms.

In practice, the greatest ideas combine both strategies: they use a moderate
number of compositional layers, each containing atoms of significant weight.
This is the ``sweet spot'' of intellectual production --- deep enough to be
non-trivial, compressed enough to be powerful.
\end{interpretation}

\bigskip

This chapter has established that depth and norm are independent measures of
complexity, explored their interaction under composition, and connected this
independence to the distinction between syntactic complexity and semantic
weight. In Chapter~\ref{ch:orthogonality}, we turn to the structure of ideas
that have zero resonance --- orthogonal ideas --- and discover the rich
geometry of unrelated thoughts.
